\begin{enunciado}{\ejercicio}
  \begin{enumerate}[label=\roman*)]
    \item Hallar todas la raíces racionales de
          \begin{enumerate}[label=(\alph*)]
            \item $2X^5 + 3X^4 + 2X^3 - X$,
            \item $X^5 - \frac{1}{2} X^4 - 2X^3 + \frac{1}{2} X^2 - \frac{7}{2}X -3$,
          \end{enumerate}
    \item Probar que $X^4 + 2X^3 - 3X^2 - 2$ no tiene raíces racionales.
  \end{enumerate}
\end{enunciado}

\begin{enumerate}[label=\roman*)]
  \item
        \begin{enumerate}[label=(\alph*)]
          \item Sale inmediato una raíz $f(r = 0) = 0$:
                $$
                  f(X) = 2X^5 + 3X^4 + 2X^3 - X = X \cdot (\ub{2X^4 + 3 X^3 + 2X^2 - 1}{g})
                $$
                Las posibles raíces de $g$ por \textit{lemma de Gauss}:
                $$
                  \set{\pm 1, \pm \frac{1}{2}}
                $$
                Probando se llega a que el conjunto de raíces $r \en \racionales$ de $f$:
                $$
                  \cajaResultado{
                    r \en \set{0, -1, \frac{1}{2}}
                  }
                $$
                Dado que no salieron más potenciales raíces del \textit{lemma de Gauss} el polinomio no tiene más raíces en $\racionales$

          \item
                Para poder usar el \textit{lemma de Gauss} hago $2f$ para tener coeficientes enteros:
                $$
                  F = 2f = 2X^5 - X^4 - 4X^3 + X^2 - 7X - 6
                $$
                Posibles raíces estarán en:
                $$
                  \big\{
                  \pm6, \pm 3, \pm 2, \pm 1, \pm \frac{3}{2}, \pm \frac{1}{2}
                  \big\}
                $$
                Probando como cualquier persona \textcolor{black!10}{sin amigos, novia, ni perro que me ladre}, encuentro 2 raíces:
                $$
                  \cajaResultado{
                    r \en \set{-\frac{3}{2}, 2}
                  }
                $$

        \end{enumerate}
  \item De tener raíces racionales estas deberían, por \textit{lemma de Gauss} estar en:
        $$
          \set{\pm2, \pm1}
        $$
        Probando esos valores ninguno es raíz
\end{enumerate}

\begin{aportes}
  \item \aporte{\dirRepo}{naD GarRaz \github}
  \item \aporte{https://github.com/ivdou}{Ivan Doumerc \github}
\end{aportes}
